\chapter{Generating Functions}

In mathematics, we have the discrete world such as natural numbers, graphs. And also the continuous world such as analysis. 
Luckily we are going to introduce a bridge between the two worlds, which is generating function. 

\section{Ordinary Generating Functions}
\begin{definition}
    Given a sequence of real numbers $a_n(n \geq 0)$. 
    We let the function $F(x)=\sum_{n=0}^{\infty}a_nx^n$ denote the \textbf{ordinary generating function} associated to $a_n$.
\end{definition}
\textbf{Remark}
\begin{itemize}
    \item The above function that we defined may not be convergent for any x.
    \item We will just use these functions as a book keeping system.
\end{itemize}

\begin{theorem}[The geometric series]
    Consider the sequence $a_n=1$ for all integers $n \geq 0$. 
    The associated generating function is $F(x)=\frac{1}{1-x}$,
\end{theorem}
\begin{proof}
    It suffices to show $(1-x)F(x)=1$.
    \begin{equation}
        \begin{split}
            (1-x)F(x) &= (1-x)\sum_{n=0}^{\infty}x^n\\
            &= \sum_{n=0}^{\infty}x^n - x \sum_{n=0}^{\infty}x^n\\
            &= \sum_{n=0}^{\infty}x^n - \sum_{n=0}^{\infty}x^{n+1}\\
            &= \sum_{n=0}^{\infty}x^n - \sum_{n=1}^{\infty}x^n\\
            &= 1
        \end{split}
    \end{equation}
\end{proof}

\begin{theorem}[The finite geometric series]
    Consider the sequence 
    \begin{equation}
        \begin{cases}
            a_n=1& \text{for all integers $0 \leq n \leq k$}\\
            a_n=0& \text{for all integers $n>k$}
        \end{cases}
    \end{equation}
    Then the associated generating function is $F(x)=\frac{1-x^{k+1}}{1-x}$
\end{theorem}
\textbf{Remark}
\begin{itemize}
    \item Writing in this way mimics inclusion-exclusion.
\end{itemize}
\begin{proof}
    \begin{equation}
        \begin{split}
            (1-x)F(x) &= (1-x)\sum_{n=0}^{k}x^n\\
            &= \sum_{n=0}^{k}x^n - x \sum_{n=0}^{k}x^n\\
            &= \sum_{n=0}^{k}x^n - \sum_{n=0}^{k}x^{n+1}\\
            &= \sum_{n=0}^{k}x^n - \sum_{n=1}^{k+1}x^n\\
            &= 1-x^{k+1}
        \end{split}
    \end{equation}
\end{proof}

\noindent \textbf{Questions} Find the generating functions.
\begin{enumerate}
    \item $\forall n \geq 0. a_n=2$
    \item $\forall n \geq 0. a_n=2^n$
    \item $a_n=1$ for $n \geq 5$ and 0 otherwise.
    \item $a_n=1$ for all $n \geq 0$ and even. $a_n=0$ otherwise.
    \item $\forall n \geq 0. a_n=2^n+1$
\end{enumerate}
\textbf{Answers}
\begin{enumerate}
    \item $F(x)=\sum_{n=0}^{\infty}a_nx^n=\sum_{n=0}^{\infty}2x^n=2\sum_{n=0}^{\infty}x^n=\frac{2}{1-x}$
    \item $F(x)=\sum_{n=0}^{\infty}a_nx^n=\sum_{n=0}^{\infty}2^nx^n=\sum_{n=0}^{\infty}(2x)^n=\frac{1}{1-2x}$
    \item define $c_n, b_n$ as below: 
    \begin{equation}
        \forall n \geq 0. c_n=1
    \end{equation} 
    \begin{equation}
        b_n = 
        \begin{cases}
            1& \text{$n<5$}\\
            0& \text{$n \geq 5$}
        \end{cases}
    \end{equation} 
    The we have $a_n=c_n-b_n$. 
    \begin{equation}
        \begin{split}
            F(x) &= \sum_{n=0}^{\infty}(c_n-b_n)x^n\\
            &= \sum_{n=0}^{\infty}c_nx^n - \sum_{n=0}^{\infty}b_nx^n\\
            &= \sum_{n=0}^{\infty}x^n - \sum_{n=0}^{4}x^n\\
            &= \frac{1}{1-x}- \frac{1-x^5}{1-x}\\
            &= \frac{x^5}{1-x}
        \end{split}
    \end{equation}
    \item \begin{equation}
        \begin{split}
            F(x) &= \sum_{n=0}^{\infty}a_nx^n\\
            &= \sum_{n=0}^{\infty}x^{2n}\\
            &= \frac{1}{1-x^2}
        \end{split}
    \end{equation}
    \item We will use linearity to calculate this.
\end{enumerate}

\begin{theorem}[Linearity]
    Take two sequences $a_n$ and $b_n$. 
    Let $F(x)$ and $G(x)$ be the associated (ordinary) generating functions respectively. 
    Let $c_n=a_n+b_n$. Then the (ordinary) generating function associated to $c_n$ is the function $F(x)+G(x)$.
\end{theorem}
\begin{proof}
    \begin{equation}
        \begin{split}
            F(x)+G(x) &= \sum_{n=0}^{\infty}a_nx^n + \sum_{n=0}^{\infty}b_nx^n\\
            &= \sum_{n=0}^{\infty}(a_nx^n+b_nx^n)\\
            &= \sum_{n=0}^{\infty}(a_n+b_n)x^n\\
            &= \sum_{n=0}^{\infty}c_nx^n\\
        \end{split}
    \end{equation}
\end{proof}

\begin{theorem}[Derivatives]
    Let $a_n$ be a sequence and let $b_{n-1}=na_n$ for $n \geq 1$. 
    If $F(x)$ is the generating function associated to $a_n$. Then $F'(x)$ is the generating function associated to $b_n$.
\end{theorem}
\begin{proof}
    Since we have $F(x)$ and we know that the derivative is linear, we can take it term by term.
    \begin{equation}
        \begin{split}
            F'(x) &= \frac{dF(x)}{dx}\\
            &= \sum_{n=0}^{\infty}na_nx^{n-1}\\
            &= \sum_{n=1}^{\infty}na_nx^{n-1}\\
            &= \sum_{n=1}^{\infty}b_{n-1}x^{n-1}\\
            &= \sum_{n=0}^{\infty}b_nx^n
        \end{split}
    \end{equation}
\end{proof}

\begin{theorem}[Multiplication]
    Let $F(x)=\sum_{n=0}^{\infty}a_nx^n$ and $G(x)=\sum_{n=0}^{\infty}b_nx^n$. 
    Then $$H(x)=F(x)G(x)= \sum_{n=0}^{\infty}c_nx^n$$ where 
    $$c_n=\sum_{m+i=n}a_mb_i=\sum_{k=0}^na_{n-k}b_k$$
\end{theorem}

\begin{theorem}[Multiple products]
    Let $F_1(x), \cdots ,F_m(x)$ be generating functions with $F_i(x)=\sum_{n=0}^{\infty}a_{n,i}x^n$. 
    Then $$F_1(x) \times \cdots \times F_m(x)=\sum_{n=0}^{\infty}c_nx^n$$ 
    where $$c_n=\sum_{e_1+ \cdots +e_m=n}a_{e_1,1}\times \cdots \times a_{e_m,m}$$
\end{theorem}

\noindent \textbf{Computation practice}
\begin{enumerate}
    \item $\frac{1}{(1-x)^2}$
    \item $\frac{1}{(1-x)^3}$
    \item $e^x$
    \item $\frac{e^x}{1+x}$
    \item $\frac{1+x+x^2}{(1-x)^{50}}$
    \item Find the number of integer solutions to the equation $e_1+e_2+e_3=25$ with $e_i \geq 0$ for all $i \in \{1,2,3\}$. $e_1$ even. $e_3 \leq 2$.
\end{enumerate}
\textbf{Answers}
\begin{enumerate}
    \item \begin{equation}
        \begin{split}
            \frac{1}{(1-x)^2} &= \frac{1}{1-x}\frac{1}{1-x}\\
            &= (1+x+x^2+x^3+ \cdots)(1+x+x^2+x^3+ \cdots)\\
            &= \sum_{n=0}^{\infty}c_nx^n
        \end{split}
    \end{equation}
    We can now use the Multiplication rule and stars and bars to calculate $c_n$.
    $$c_n=\sum_{e_1+e_2=n}1 \times 1={n+(2-1) \choose 2-1}={n+1 \choose 1}=n+1$$
    \item Similar to the last question $$c_n = \sum_{e_1+e_2+e_3}1 \times 1 \times 1={n+2 \choose 2}$$
    \item Using taylor expansion $$e^x=\sum_{n=0}^{\infty}\frac{1}{n!}x^n$$
    \item \begin{equation}
        \begin{split}
            \frac{e^x}{1+x} &= (\sum_{n=0}^{\infty}\frac{1}{n!}x^n)\times (\sum_{n=0}^{\infty}(-x)^n)\\
            &= \sum_{n=0}^{\infty}c_nx^n
        \end{split}
    \end{equation}
    where $$c_n=\sum_{e_1+e_2=n}\frac{1}{e_1!}(-1)^{e_2}=\sum_{k=0}^n\frac{1}{k!}(-1)^{n-k}$$
    \item \begin{equation}
        \begin{split}
            \frac{1+x+x^2}{(1-x)^{50}} &= \frac{1}{(1-x)^{50}} + \frac{x}{(1-x)^{50}} + \frac{x^2}{(1-x)^{50}}\\
            &= {n+49 \choose 49} + {n+48 \choose 49} + {n+47 \choose 49}
        \end{split}
    \end{equation}
    \item Let us first compute the generationg function $F(x)=\sum_{n=0}^{\infty}a_nx^n$ associated to the solutions $\{a_n\}$. 
    $\{a_n\}$ is a sequence, where $a_n$ = number of solutions of $e_1+e_2+e_3=n$.
    \begin{itemize}
        \item Solutions of $e_1$: $S_{e_1}=\{0,2,4,6, \cdots\}$
        \item Solutions of $e_2$: $S_{e_2}=\{0,1,2,3, \cdots\}$
        \item Solutions of $e_3$: $S_{e_3}=\{0,1,2\}$
    \end{itemize}
    Then we only need to show $a_{25}$. i.e. choose one of $S_{e_1},S_{e_2},S_{e_3}$ to make 25. 
    This is equivalent to find the coefficient of $x^{25}$ for the function below.
    $$(x^0+x^2+x^4+ \cdots)(x^0+x^1+x^2+\cdots)(x^0+x^1+x^2)=\frac{1}{1-x^2}\frac{1}{1-x}\frac{1-x^3}{1-x}$$
\end{enumerate}

\begin{definition}
    Given a power series $f(x)=\sum_{n=0}^{\infty}a_nx^n$, 
    we often write the shorthand $[x^n](f(x))=a_n$ instead of saying the n-th coefficient.
\end{definition}
\textbf{Remark}
\begin{itemize}
    \item We can view this notation as a function where the input is the function $f$ and the output is the coefficient.
    \item $[x^n](f(x)+g(x))=[x^n](f(x))+[x^n](g(x))$
\end{itemize}

\begin{theorem}
    $\forall p \in \N \cup\{0\}. [x^n]((1+x)^p)={p \choose n}$
\end{theorem}
\begin{proof}
    Recall the Binomial theorem and use the definition above.
\end{proof}

\section{Newton's Binomial Theorem}

\begin{theorem}[Newton's Binomial Theorem]
    For every \underline{non-zero real number} p. 
    We have the power series expansion of $(1+x)^p$
    $$\sum_{k=0}^{\infty}\frac{f^{(k)}(0)}{k!}x^k$$
    We let ${p \choose n}$ denote the unique n-th coefficient in the power series.
\end{theorem}
\textbf{Remark}
\begin{itemize}
    \item Normally, one can also define $P(p,n)={p \choose n}n!$
\end{itemize}

\begin{lemma}
    For all real p for which it is defined. we have ${p \choose 0}=1$.
\end{lemma}
\begin{proof}
    ${p \choose 0}=\frac{f^{(0)}(0)}{0!}=\frac{f(0)}{1}=(1+0)^p=1$
\end{proof}

\begin{lemma}
    For all real p for which it is defined and $n \in \N$.
    $${p \choose n}=\frac{p}{n}{p-1 \choose n-1}$$
\end{lemma}
\begin{proof}:
    \begin{itemize}
        \item For $p \in \N$, it is trivial. $${p \choose n}=\frac{p!}{n!(p-n)!}=\frac{p}{n}\frac{(p-1)!}{(n-1)!(p-n)!}$$
        \item We focus on $p \notin \N$. By Newton's Binomial Theorem we have $$(1+x)^p=\sum_{k=0}^{\infty}{p \choose n}x^n$$
        Let us take the derivative for both sides.
        \begin{equation}
            \begin{split}
                p(1+x)^{p-1} &= \sum_{n=0}^{\infty}{p \choose n}nx^{n-1}\\
                (1+x)^{p-1} &= \sum_{n=0}^{\infty}\frac{n}{p}{p \choose n}x^{n-1}
            \end{split}
        \end{equation}
        Find the coefficient for $x^{n-1}$ for both sides
        \begin{equation}
            \begin{split}
                [x^{n-1}]((1+x)^{p-1}) &= {p-1 \choose n-1}\\
                [x^{n-1}](\sum_{n=0}^{\infty}\frac{n}{p}{p \choose n}x^{n-1}) &= \frac{n}{p}{p \choose n}
            \end{split}
        \end{equation}
        Hence we have $${p-1 \choose n-1}= \frac{n}{p}{p \choose n}$$
        Therefore, $${p \choose n}= \frac{p}{n}{p-1 \choose n-1}$$
    \end{itemize}
\end{proof}

\noindent \textbf{Examples}
Using the above theorems and lemmas, we can efficiently compute some fractinal binomial coefficients.
\begin{enumerate}
    \item ${-5 \choose 3}$
    \item ${\frac{1}{2} \choose 3}$
\end{enumerate}
\textbf{Answers}
\begin{enumerate}
    \item \begin{equation}
        \begin{split}
            {-5 \choose 3} &= \frac{-5}{3}\cdot{-6 \choose 2}\\
            &= \frac{-5}{3}\cdot\frac{-6}{2}\cdot{-7 \choose 1}\\
            &= \frac{-5}{3}\cdot\frac{-6}{2}\cdot\frac{-7}{1}{-8 \choose 0}
        \end{split}
    \end{equation}
    \item \begin{equation}
        \begin{split}
            {\frac{1}{2} \choose 3} &= \frac{\frac{1}{2}}{3} \cdot {-\frac{1}{2} \choose 2}\\
            &= \frac{\frac{1}{2}}{3}\cdot\frac{-\frac{1}{2}}{2}\cdot{-\frac{3}{2} \choose 1}\\
            &= \frac{\frac{1}{2}}{3}\cdot\frac{-\frac{1}{2}}{2}\cdot\frac{-\frac{3}{2}}{1}\cdot{-\frac{5}{2} \choose 0}
        \end{split}
    \end{equation}
\end{enumerate}

\begin{theorem}
    For $p<0$ $${p \choose n}= (-1)^n{-p+n-1 \choose n}$$
\end{theorem}
\begin{proof}
    By induction.
\end{proof}

\begin{corollary}
    For $p>0$ $$\frac{1}{(1+x)^p}=\sum_{n=0}^{\infty}(-1)^n{p+n-1 \choose n}x^n$$
\end{corollary}

\begin{lemma}
    $$(1-4x)^{-\frac{1}{2}}=\sum_{n=0}^{\infty}{2n \choose n}x^n$$
\end{lemma}
\begin{proof}
    By Newton's Binomial Theorem we have: $$(1-4x)^{-\frac{1}{2}}=\sum_{n=0}^{\infty}{-\frac{1}{2} \choose n}(-2x)^n$$
    So it suffices to show that $${-\frac{1}{2} \choose n}(-4)^n={2n \choose n}$$
    We finish the proof by showing:
    \begin{equation}
        \begin{split}
            {-\frac{1}{2} \choose n}(-4)^n &= (-1)^n \cdot 2^{2n} \cdot {-\frac{1}{2} \choose n}\\
            &= (-1)^n \cdot 2^{2n} \cdot \frac{-\frac{1}{2}}{n}{-\frac{1}{2}-1 \choose n-1}\\
            &= (-1)^n \cdot 2^{2n} \cdot \frac{-\frac{1}{2}}{n} \cdot \frac{-\frac{1}{2}-1}{n-1} \cdot {-\frac{1}{2}-1 \choose n-1}\\
            &= (-1)^n \cdot 2^{2n} \cdot \frac{-\frac{1}{2}(-\frac{1}{2}-1) \cdots (-\frac{1}{2}-(n-1))}{n(n-1) \cdots 1}\\
            &= (-1)^n \cdot 2^{2n} \cdot \frac{(-\frac{1}{2})(-\frac{3}{2}) \cdots (-\frac{1-2n}{2})}{n!}\\
            &= (-1)^n \cdot 2^{2n} \cdot \frac{(-1)(-3) \cdots (1-2n)}{2^n \cdot n!}\\
            &= (-1)^n \cdot 2^{2n} \cdot \frac{(-1)^n(1 \times 3 \times \cdots \times (2n-1))}{2^n \cdot n!}\\
            &= (-1)^n \cdot 2^{2n} \cdot \frac{(-1)^n}{2^n} \cdot \frac{1 \times 2 \times3 \times 4 \times \cdots \times (2n-1) \times 2n}{n!} \cdot \frac{1}{2 \times 4 \times 6 \times \cdots \times 2n}\\
            &= 2^n \cdot \frac{(2n)!}{n!} \cdot \frac{1}{2^n \cdot n!}\\
            &= \frac{(2n)!}{n!(2n-n)!}\\
            &= {2n \choose n}
        \end{split}
    \end{equation}
\end{proof}

\section{Partition of integers}

We have learnt some properties of integers from elimentary school. This times we will use the ordinary generating functions as a powerful method to find some interesting facts.

\begin{definition}
    A \textbf{partition of an integer $n$} is a multiset $P$. 
    (a set that is allowed to have the same element multiple times) 
    with the property that $$\sum_{i \in P}i=n \text{ with } i>0$$
\end{definition}
\textbf{Remark}
\begin{itemize}
    \item The key idea here is a partition is a splitting of an integer into pieces, 
    where order does not matter, contrary to distribution problems.
    \item $\{1,2,1,1\}=\{2,1,1,1\}$ sometimes we like decending order.
\end{itemize}
\textbf{Example:} partitions of the integer 5:
$$\{5\},\{4,1\},\{3,2\},\{2,2,1\},\{3,1,1\},\{2,1,1,1\},\{1,1,1,1,1\}$$

There is a better way to view partitions:
\begin{definition}
    A \textbf{partition of an integer $n$} is an infinite string of integers $e_i \geq 0$ 
    indexed by $i \in \N$ with the property that $$\sum_{i=1}^{\infty}i \cdot e_i=n$$
\end{definition}
\textbf{Remark:}
\begin{itemize}
    \item These strings count the frequency of a piece of size $e_i$. i.e. $i$ repeated $e_i$ times.
    \item This is more formal than multisets and helps you see the next theorems more clearly.
    \item We can see some similar constructions in some diagonalization proofs.
\end{itemize}
\textbf{Example:} Still consider the partition of the integer 5:
\begin{itemize}
    \item $\{5\}: 00001000 \cdots $
    \item $\{4,1\}: 10010000 \cdots $
    \item $\{3,1,1\}: 20100000 \cdots $
\end{itemize}

\begin{definition}
    We use $p_n$ to denote the number of partitions of the integer n.
\end{definition}

\begin{theorem}
    Let $P(x)=\sum_{n=0}^{\infty}p_nx^n$. 
    We can then express $P(x)$ as the infinite product $$P(x)=\prod_{n=1}^{\infty}\frac{1}{1-x^n}$$
\end{theorem}
\begin{proof}
    Consider the infinite product 
    \begin{equation}
        \begin{split}
            \prod_{n=1}^{\infty}\frac{1}{1-x^n} &= \frac{1}{1-x} \times \frac{1}{1-x^2} \times \frac{1}{1-x^3} \times \cdots\\
            &= (\sum_{k=0}^{\infty}x^n)\times(\sum_{k=0}^{\infty}(x^2)^n)\times(\sum_{k=0}^{\infty}(x^3)^n) \times \cdots\\
            &= \sum_{k=0}^{\infty}c_nx^n
        \end{split}
    \end{equation}
    We let  $c_n$ denotes the coefficient of $x^n$, and it can be computed associated
    \begin{equation}
        \begin{split}
            c_n &= \sum_{e_1+2e_2+3e_3+ \cdots =n}1 \times 1 \times 1 \times \cdots\\
            &= \sum_{e_1+2e_2+3e_3+ \cdots =n}1\\
            &= p_n
        \end{split}
    \end{equation}
\end{proof}

\begin{definition}
    We say a partition of $n$.(i.e. $\sum_{i=1}^{\infty}i \cdot e_i=n$) has \textbf{distinct parts} if $\forall i \in \N.e_i \in \{0,1\}$
\end{definition}

\begin{definition}
    Let $d_n$ be the number of partitions of $n$ with distinct parts.
\end{definition}

\begin{definition}
    We say a partition of $n$ only has \textbf{odd parts} if $e_i=0$ whenever $i$ is even. Equivalently, $n=\sum_{i=1}^{\infty}(2i-1)e_{2i-1}$
\end{definition}

\begin{definition}
    Let $o_n$ be the number of partitions of $n$ with odd parts.
\end{definition}

\noindent\textbf{Example} What is $o_5$?
\begin{enumerate}
    \item $\{5\}$: True
    \item $\{3,2\}$: False, has 2.
    \item $\{4,1\}$: False, has 4.
    \item $\{3,1,1\}$: True
    \item $\{2,2,1\}$: False, has 2.
    \item $\{2,1,1,1\}$: False, has 2.
    \item $\{1,1,1,1,1\}$: True.
\end{enumerate}
So $o_5=3$.

\begin{theorem}
    Let $$O(x)=\sum_{n=0}^{\infty}o_nx^n$$
    We can then express $O(x)$ as the infinite product $$O(x)=\prod_{k=1}^{\infty}\frac{1}{1-x^{2k-1}}$$
\end{theorem}

\begin{theorem}
    Let $$D(x)=\sum_{n=0}^{\infty}d_nx^n$$
    We can then express $D(x)$ as the infinite product $$D(x)=\prod_{k=1}^{\infty}(1+x^k)$$
\end{theorem}
\begin{proof}
    Consider the equation $$e_1+2e_2+3e_3+4e_4+ \cdots=n \text{ where } e_i \in {0,1}$$
    Let $S(n)$ denotes the number of solutions for the equation. \\
    Then the sequence $\{S(n)\}$ can be associated with a generating function $$\sum_{n=0}^{\infty}S(n)x^n=\sum_{n=0}^{\infty}d_nx^n$$
    We notice that $$S(n)=[x^n]((1+x)(1+x^2)(1+x^3)\cdots)$$
    This is because we can associate $e_i$ with the factor $(1+x^{e_i})$.
    \begin{itemize}
        \item if $e_i=0$, then the factor will contribute $x^0$, which is 1.
        \item if $e_i \neq 0$, then the factor will contribute $x^{e_i}$
    \end{itemize}
    Hence $$D(x)=\sum_{n=0}^{\infty}d_nx^n=(1+x)(1+x^2)(1+x^3)\cdots=\prod_{k=1}^{\infty}(1+x^k)$$
\end{proof}

\noindent The last theorem is about \underline{the connection between distinc and odd}.
\begin{theorem}
    $$D(x)=O(x)$$
    Consequently, the number of partitions of n into odd parts is equal to the number of partitions of n into distinct parts.
\end{theorem}
\textbf{Remark}
\begin{itemize}
    \item We do not know the bijective correspondence between partitions into odd parts and partitions into distinct parts.
    \item We know they are equal because they have the same generating functions.
\end{itemize}
\begin{proof}
    We already showed that $$D(x)=\prod_{k=1}^{\infty}(1+x^k)$$
    Now let's consider $O(x)$
    \begin{equation}
        \begin{split}
            O(x) &= \prod_{n=1}^{\infty}\frac{1}{1-x^{2n-1}}\\
            &= \prod_{n=1}^{\infty}\frac{1}{1-x^{2n-1}} \cdot \prod_{n=1}{\infty}\frac{1-x^{2n}}{1-x^{2n}}\\
            &= \frac{\prod_{n=1}^{\infty}(1-x^{2n})}{\prod_{n=1}^{\infty}(1-x^n)}\\
            &= \prod_{n=1}^{\infty}\frac{(1-x^n)(1+x^n)}{(1-x^n)}\\
            &= \prod_{n=1}^{\infty}(1+x^n)\\
            &= D(X)
        \end{split}
    \end{equation}
\end{proof}


\section{Exponential Generating Functions}

\begin{definition}
    Given a sequnce $\{a_n\}$ where $n \in \N \cup \{0\}$. We call $$f(x)=\sum_{n=0}^{\infty}\frac{a_n}{n!}x^n$$ \textbf{the exponential generating function} associated to $a_n$.
\end{definition}
\textbf{Remark}
\begin{itemize}
    \item Why divide by $n!$ is because in some way, we should think of canceling out possible permutations.
    \item We can have a wider range of possible sequences that we can use. Since we divide $a_n$ by $n!$, this is now able to converge more.
    \item $f^{(n)}(0)=a_n$ this is nice.
    \item $\forall n \in \N.a_n=1$ then $f(x)=e^x$
\end{itemize}

\noindent \textbf{Examples} Find the exponential generating functions.
\begin{enumerate}
    \item $\forall n \geq 0.a_n=2$
    \item $\forall n \geq 0.a_n=2^n$
    \item $\forall n \geq 0.a_n=(-1)^n$
    \item $a_n=1$ for all $n \leq 5$ and $a_n=0$ otherwise.
    \item $\forall n \geq 0.a_n=n!$
    \item $a_n=1$ when n is even and 0 otherwise.
\end{enumerate}
\textbf{Answer}
\begin{enumerate}
    \item $$
            f(x) = \sum_{n=0}^{\infty}\frac{2}{n!}x^n
            = 2 \sum_{n=0}^{\infty}\frac{x^n}{n!}
            = 2e^x
    $$
    \item $$
        f(x) = \sum_{n=0}^{\infty}\frac{2^n}{n!}x^n
        = \sum_{n=0}^{\infty}\frac{(2x)^n}{n!}
        = e^{2x}
    $$
    \item $$
        f(x) = \sum_{n=0}^{\infty}\frac{(-1)^n}{n!}x^n
        = \sum_{n=0}^{\infty}\frac{(-x)^n}{n!}
        = e^{-x}
    $$
    \item $$f(x)=1+x+\frac{x^2}{2}+\frac{x^3}{6}+\frac{x^4}{24}+\frac{x^5}{120}$$
    \item $$
    f(x) = \sum_{n=0}^{\infty}\frac{a_n}{n!}x^n
    = \sum_{n=0}^{\infty}x^n
    = \frac{1}{1-x}
    $$
    \item We notice that $$\sum_{n=0}^{\infty}\frac{x^n}{n!}+\sum_{n=0}^{\infty}\frac{(-1)^nx^n}{n!}=2\sum_{n=0}^{\infty}\frac{a_n}{n!}x^n$$
    Hence we have $$e^x+e^{-x}=2f(x)$$
    Therefore we get $$f(x)=\frac{e^x+e^{-x}}{2}=cosh(x)$$
\end{enumerate}

\begin{theorem}[Linearity]
    Let $a_n$ and $b_n$ be sequences, let $f(x)$ and $g(x)$ be the corresponding exponential generating functions. 
    The exponential generating function associated to $c_n =a_n+b_n$ is $f(x)+g(x)$.
\end{theorem}

\begin{theorem}[Multiplication]
    Let $$f(X)=\sum_{n=0}^{\infty}\frac{a_n}{n!}x^n$$
    $$g(X)=\sum_{n=0}^{\infty}\frac{b_n}{n!}x^n$$
    Then the function $f(X)g(X)$ is the exponential generating function for the sequence $$c_n=\sum_{k=0}^{\infty}{n \choose k}a_kb_{n-k}$$
\end{theorem}
\begin{proof}
    \begin{equation}
        \begin{split}
            f(x)g(x) &= (\sum_{n=0}^{\infty}\frac{a_n}{n!}x^n)(\sum_{n=0}^{\infty}\frac{b_n}{n!}x^n)\\
            &= \sum_{n=0}^{\infty}(\sum_{k=0}^{n}\frac{a_k}{k!} \cdot \frac{b_{n-k}}{(n-k)!})x^n\\
            &= \sum_{n=0}^{\infty}n!(\sum_{k=0}^{n}\frac{a_k}{k!} \cdot \frac{b_{n-k}}{(n-k)!})\frac{x^n}{n!}\\
            &= \sum_{n=0}^{\infty}(\sum_{k=0}^{n}\frac{n!}{k!(n-k)!}a_kb_{n-k})\frac{x^n}{n!}\\
            &= \sum_{n=0}^{\infty}(\sum_{k=0}^{\infty}{n \choose k}a_kb_{n-k})\frac{x^n}{n!}
        \end{split}
    \end{equation}
\end{proof}

\begin{theorem}[Multiple products]
    For all $i \in \{1, \cdots ,k\}$, let $$f_i(x)=\sum_{n=0}^{\infty}\frac{a_{e_i,i}}{n!}x^n$$
    Then $$\prod_{i=1}^{k}f_i(x)=\sum_{n=0}^{\infty}\frac{c_n}{n!}x^n$$
    where $$c_n=\sum_{e_1+ \cdots +e_k=n}\frac{n!}{e_1! \times \cdots \times e_k!}a_{e_1,1}\times \cdots \times a_{e_k,k}$$
\end{theorem}

\noindent \textbf{Example}: Compute the number of binary strings there are with length n.\\
\textbf{Answer}: It is very confusing for students that neither the book nor the professor explained why we determine such exponential generating function. I will try my best to explain.
\begin{itemize}
    \item We already know the answer is $${n \choose 0}+{n \choose 1}+ \cdots +{n \choose n}=\sum_{k=0}^{n}{n \choose k}$$ because we can consider the number of 0's. Hence the coefficient we need is $$c_n=\sum_{k=0}^n {n \choose k}$$
    \item What is the difference of considering the coefficient of $x^n$ and the coefficient of $\frac{x^n}{n!}$?\\
    We let $A_n=\frac{a_n}{n!}$, $B_n=\frac{b_n}{n!}$, $C_n=\frac{c_n}{n!}$
    $$(\sum_{n=0}^{\infty}A_nx^n)(\sum_{n=0}^{\infty}B_nx^n)=\sum_{n=0}^{\infty}C_nx^n \Longrightarrow C_n=\sum_{e_1+e_2=n}A_{e_1}B_{e_2}$$
    $$(\sum_{n=0}^{\infty}\frac{a_n}{n!}x^n)(\sum_{n=0}^{\infty}\frac{b_n}{n!}x^n)=\sum_{n=0}^{\infty}\frac{c_n}{n!}x^n \Longrightarrow c_n=\sum_{e_1+e_2=n}\frac{n!}{e_1! \cdot e_2!}$$
    So considering the coefficient of $\frac{x^n}{n!}$ means we considering the order of every spot and then eliminate the order in 0's and 1's.
    Therefore $$c_n=\sum_{e_1+e_2=n}\frac{n!}{e_1! \cdot e_2!}=\sum_{k=0}^{n}{n \choose k}$$
    \item Let $e_0$ denotes the number of 0's.\\
    Let $e_1$ denotes the number of 1's.\\
    We should consider 
    \begin{equation}
        \begin{cases}
            e_0+e_1=n \text{ with } e_0,e_1 \geq 0\\
            \text{the distribution of 0's and 1's}
        \end{cases}
    \end{equation}
    Since we need to consider the distributions, as talked above, we need to use the exponential generating functions.\\
    We determine the exponential generating function for digit 0:
    $$1+\frac{x}{1!}+\frac{x^2}{2!}+\frac{x^3}{3!}+ \cdots$$
    We determine the exponential generating function for digit 1:
    $$1+\frac{x}{1!}+\frac{x^2}{2!}+\frac{x^3}{3!}+ \cdots$$
    The exponential generating function for our answer is:
    $$(1+\frac{x}{1!}+\frac{x^2}{2!}+\frac{x^3}{3!}+ \cdots)(1+\frac{x}{1!}+\frac{x^2}{2!}+\frac{x^3}{3!}+ \cdots)=e^{2x}$$
\end{itemize}

\noindent \textbf{Example}: Compute the exponential generating function associated to ternary strings with at least five zeros.\\
(Hint: mentally piggyback off the idea of a distribution $e_0+e_1+e_2=n$ where order matters and $e_0 \geq 5$. Convince yourself that the distribution of the coefficients are counting what you want them to.)\\
\textbf{Answer}:\\
Let $e_0$ denotes the number of 0's.\\
Let $e_1$ denotes the number of 1's.\\
Let $e_2$ denotes the number of 2's.\\
We should consider
\begin{equation}
    \begin{cases}
        e_0+e_1+e_2=n \text{ with } e_1,e_2 \geq 0 \text{ and }e_0 \geq 5\\
        \text{the distribution of 0's, 1's and 2's}
    \end{cases}
\end{equation}
We determine the generating function for digit 0:
$$\frac{x^5}{5!}+\frac{x^6}{6!}+\frac{x^7}{7!}+ \cdots$$
We determine the generating function for digit 1:
$$1+\frac{x}{1!}+\frac{x^2}{2!}+\frac{x^3}{3!}+ \cdots$$
We determine the generating function for digit 2:
$$1+\frac{x}{1!}+\frac{x^2}{2!}+\frac{x^3}{3!}+ \cdots$$
The answer we want:
$$(\frac{x^5}{5!}+\frac{x^6}{6!}+\frac{x^7}{7!}+ \cdots)(1+\frac{x}{1!}+\frac{x^2}{2!}+\frac{x^3}{3!}+ \cdots)^2=(e^x-1-\frac{x^2}{2!}-\frac{x^3}{3!}-\frac{x^4}{4!})e^{2x}$$

\noindent \textbf{Example}: Compute the generating function associated to case sensitive alpha-numeric passwords containing at least one capital and one number.\\
\textbf{Answer}:\\
Let $e_a, \cdots ,e_z$ denotes the number of a's to the number of z's respectively.\\
Let $e_A, \cdots ,e_Z$ denotes the number of A's to the number of Z's respectively.\\
Let $e_0, \cdots ,e_9$ denotes the number of 0's to the number of 9's respectively.\\
We should consider
\begin{equation}
    \begin{cases}
        (e_a+ \cdots +e_z)+(e_A+ \cdots +e_Z)+(e_0+ \cdots +e_9)=n\\
        e_a, \cdots ,e_z \geq 0\\
        e_A, \cdots ,e_Z \geq 0 \text{ AND not all of } e_{\alpha}=0 (\alpha \in \{A, \cdots ,Z\})\\
        e_0, \cdots ,e_9 \geq 0 \text{ AND not all of } e_{\beta}=0 (\beta \in \{0, \cdots ,9\})\\
        \text{the order distribution}
    \end{cases}
\end{equation}
All the generating functions for $e_a$ to $e_z$ are:
$$1+\frac{x}{1!}+\frac{x^2}{2!}+\frac{x^3}{3!}+ \cdots$$
Hence the generating function for the part $(e_a+ \cdots +e_z)$ is:
$$(1+\frac{x}{1!}+\frac{x^2}{2!}+\frac{x^3}{3!}+ \cdots)^{26}=e^{26x}$$
If we don't consider not all of  $e_{\alpha}=0 (\alpha \in \{A, \cdots ,Z\})$, then All the generating functions for $e_A$ to $e_Z$ are:
$$1+\frac{x}{1!}+\frac{x^2}{2!}+\frac{x^3}{3!}+ \cdots$$
In this situation, the exponential generating function for the part $(e_A+ \cdots +e_Z)$ is also:
$$(1+\frac{x}{1!}+\frac{x^2}{2!}+\frac{x^3}{3!}+ \cdots)^{26}=e^{26x}$$
But if we consider that condition, we should minus the case that all uppercase letters are not chosen, hence the coefficient for $x^0$ in the exponential generating function for the part $(e_A+ \cdots +e_Z)$ should minus 1.
Therefor, the exponential generating function for the part $(e_A+ \cdots +e_Z)$ is:
$$(1+\frac{x}{1!}+\frac{x^2}{2!}+\frac{x^3}{3!}+ \cdots)^{26}-x^0=e^{26x}-1$$
Similarly, the exponential generating function for the part $(e_0+ \cdots +e_9)$ is:
$$(1+\frac{x}{1!}+\frac{x^2}{2!}+\frac{x^3}{3!}+ \cdots)^{10}-x^0=e^{10x}-1$$
Therefore, the answer is:
$$e^{26x}\times (e^{26x}-1)\times (e^{10x}-1)=e^{62x}-e^{36x}-e^{52x}+e^{26x}$$
Also, we can calculate explicitly that $$a_n=62^n-36^n-52^n+26^n$$

\noindent \textbf{Example}: Compute the generating function associated to ternary strings with even number of 0's and odd number of 1's.\\
\textbf{Answer}:\\
Let $e_0$ denotes the number of 0's.\\
Let $e_1$ denotes the number of 1's.\\
Let $e_2$ denotes the number of 2's.\\
We should consider
\begin{equation}
    \begin{cases}
        e_0+e_1+e_2=n\\
        e_0 \geq 0 \text{ and even}\\
        e_1 \geq 0 \text{ and odd}\\
        e_2 \geq 0\\
        \text{the distribution of 0's, 1's and 2's}
    \end{cases}
\end{equation}
We determine the generating function for digit 0:
$$1 \cdot \frac{x^0}{0!} + 0 \cdot \frac{x^1}{1!} + 1 \cdot \frac{x^2}{2!}+ 0 \cdot \frac{x^3}{3!} + \cdots=\frac{e^x+e^{-x}}{2}$$
There are two ways to determine the generating function for digit 1:
\begin{itemize}
    \item no restriction minus the even cases: $$e^x-\frac{e^x+e^{-x}}{2}=\frac{e^x-e^{-x}}{2}$$
    \item Consider the even cases: $$1 \cdot \frac{x^0}{0!} + 0 \cdot \frac{x^1}{1!} + 1 \cdot \frac{x^2}{2!}+ 0 \cdot \frac{x^3}{3!} + \cdots=\frac{e^x+e^{-x}}{2}$$
    Then taking the derivative: $$x+\frac{x^3}{3!}+\frac{x^5}{5!}+ \cdots= \frac{e^x-e^{-x}}{2}$$
\end{itemize}
We determine the generating function for digit 2:
$$1+\frac{x}{1!}+\frac{x^2}{2!}+\frac{x^3}{3!}+ \cdots=e^x$$
Therefore the answer is:
$$\frac{e^x+e^{-x}}{2} \times \frac{e^x-e^{-x}}{2} \times e^x$$